% ============================================================================
% Глава 1: Въведение
% ============================================================================

\section{ВЪВЕДЕНИЕ}

В съвременната образователна среда платформите за електронно обучение стават все по-важни за предоставянето на персонализирани учебни преживявания. Тези системи трябва да управляват сложни иерархични данни за студенти, курсове, учебни модули и напредък, които често имат променлива структура. Традиционните релационни бази данни срещат затруднения при обработката на такива гъвкави и вложени структури от данни \cite{sadalage2012nosql}.

Бизнес проблемът, който решава този проект, е необходимостта от ефективно съхранение и обработка на данни в платформа за персонализирано обучение. Системата трябва да проследява напредъка на студенти по различни курсове, включително детайлни метаданни за активността им, резултати от тестове и алтернативни методи за завършване на курсове. Това изисква гъвкава структура, която може да се адаптира към различни образователни сценарии без необходимост от чести промени в схемата на базата данни.

Основните цели на проекта включват:
\begin{enumerate}
    \item Създаване на работеща база данни в MongoDB с правилно структурирани данни \cite{mongodb_docs}
    \item Реализация на CRUD операции (създаване, четене, актуализиране, изтриване) с различни видове филтриране
    \item Разработка на агрегиращи заявки за анализ на данни
    \item Внедряване на механизъм за управление на достъпа чрез роли
    \item Създаване на REST API и уеб интерфейс за демонстрация на функционалностите
\end{enumerate}

Изборът на нерелационна база данни е оправдан от естеството на данните в образователната сфера. MongoDB предлага следните предимства пред традиционните релационни системи:
\begin{itemize}
    \item Гъвкавост в структурата на документите, позволяваща лесно добавяне на нови полета без миграции
    \item Възможност за влагане на документи (embedded documents) за представяне на йерархични структури като модули и лекции в курсове
    \item По-добра производителност при работа с големи обеми данни благодарение на хоризонталното мащабиране
    \item Подходящ модел за данни, който естествено отразява обектно-ориентирания подход в съвременните приложения
\end{itemize}

Проектът демонстрира как нерелационните бази данни могат да бъдат ефективно използвани за решаване на реални бизнес проблеми в сферата на електронното обучение, като същевременно предоставя солидна основа за бъдещо разширение и развитие на системата.
